\scanpage{116}
\mspage{57}
\margin{(174)\quad \figure{\begin{tikzpicture}[x=.5cm,y=.5cm,line width=.35pt]\draw (0,9.2)..controls (.1,7.6) and (-.1,4)..(.03,2.7)..controls (.04,2.5) and (.12,2.35)..(.35,2.25);\draw (.03,2.7)..controls (.12,1.8) and (.45,.2)..(.85,0);\end{tikzpicture}}}
\diagram{\(\unsure{(}\quad\)\begin{tikzcd}[ampersand replacement=\&, column sep=large,row sep=small]
D_{d,*} \arrow[r,"\sigma_d"] \arrow[d,equal,"\mathrm{def}"'] \& D_d \\
\displaystyle\coprod_{\substack{\deleted{\unsure{p\leq d|d'}}\\d'\in\mathbb N}}D_{d,d'}=\sigma^{-1}(D_d) \&
\end{tikzcd}}
qui fait de $D_{d,*}$ un esp. analytique sur $D_d$, \deleted{\unsure{pondéré par $d'$}} relativement ordonné (par la relation de factorisabilité
\diagram{\begin{tikzcd}[ampersand replacement=\&, column sep=small,row sep=small]
{} \& {} \arrow[dd,dotted] \\
\bullet \arrow[ru] \arrow[rd] \& \\
{} \& {}
\end{tikzcd}}
des flèches de source donnée) \inserted{``pondéré'' par $d'$,} et loc. isomorphe à $\deleted{\unsure{E}}\mathcal P(I_d)^{\deleted{\unsure{0|\circ}}}_{D_d}$ \inserted{avec sa relation d'ordre et la fonction ``cardinal''} (--- le faisceau constant d'ensembles ordonnés simpliciaux $\mathcal P(I_d)^\circ$ ---

Cela redonne donc un rev. $D^!_d$ principal de groupe $\mathfrak S_d$, permettant de reconstituer les $D_{d,d'}$ --- mais pour reconstituer $D$, il faut de plus se donner \deleted{\unsure{les}} au moins les $D_{d,d'}\unsure{\text{---}|\longrightarrow}D_{d'}$, et (c'est plus sûr!) aussi
\[
D_{\unsure{d|d^\circ},d'}\mathbin{\times}_{D_{d'}}D_{d',d''}\unsure{\text{---}|\longrightarrow}D_{\unsure{d|d^\circ},d''}
\]
(la \unsure{flèche|flèches} composition), car il n'est pas clair, avec les seules hyp. faites jusqu'à présent, que \unsure{le|les} \unsure{morphisme|morphismes}
\[
D_{d,d''}\unsure{\text{---}|\longrightarrow}D_d\times D_{d''}
\]
sont mono ?, i.e. que ``tt flèches sont déterminées par sa source et par son but''. Mais finalement si, c'est une conséquence de c). Donc
\margin{115}
